\section{Hierarchy Counterexamples}
\label{app:hierarchy-counterexamples}

This appendix records counterexamples that support the hierarchy discussion in
\Cref{sec:stability-hierarchy}.  Each example is finite, explicit, and
hand-checkable under the conventions of \Cref{sec:perturbation-semantics}.

\subsection{Delete-one does not imply insert-one}

Take metric space $\{a, b\}$ with $d(a,b) = 1$.  Let $S = ((a,0), (b,0))$ and
fix $k=1$, query $(a,1)$.  The original prediction at $a$ is $0$, so the loss
at $(a,1)$ is $1$.  Deleting either index leaves a sample whose $1$-NN
prediction at $a$ is still $0$.  Thus $\Delta_{\mathrm{del}} = 0$ for all
$i,x,y$.

Inserting $(a,1)$ at position~1 gives $((a,1), (a,0), (b,0))$, and the $1$-NN
prediction at $a$ becomes $1$, changing the loss at $(a,1)$ from $1$ to $0$.
Hence $\Delta_{\mathrm{ins}} = 1$.

\subsection{Insert-one does not imply delete-one}

Take $k=2$ with metric space $\{a,b\}$, $d(a,b)=1$.  Let $S = ((a,0), (b,1))$
and fix query $(a,1)$.  The original $2$-NN prediction at $a$ is $0$ (margin
$0$, tie resolved to $0$), so the loss at $(a,1)$ is $1$.  Deleting index $0$
leaves $((b,1))$; with $k'=1$, the $1$-NN prediction at $a$ is $1$, changing
the loss to $0$.  Hence $\Delta_{\mathrm{del}} = 1$ at
$(i=0,\;x=a,\;y=1)$.  For any insertion at any position, the $2$-NN at $a$
always includes an $(a,0)$ before any inserted point at distance $0$, so the
margin stays at $0$ and the prediction remains $0$ by tie-breaking.  Thus
$\Delta_{\mathrm{ins}} = 0$ for all insertions.

\subsection{LOO does not imply pointwise delete-one}

Let the metric space be the three-point linear metric $\{0,1,2\}$ with
$d(0,1)=d(1,2)=1$, $d(0,2)=2$.  Fix $k=1$ and
$S = ((0,0), (2,1), (0,0), (2,1))$.  Every point appears with a same-label
duplicate, so LOO at each index is $0$.

Now consider delete-one at index~0 with query $(1,0)$.  The original $1$-NN at
$1$ is a tie between $(0,0)$ (index~0) and $(2,1)$ (index~1), resolved to
$(0,0)$, predicting $0$; loss at $(1,0)$ is $0$.  After deleting index~0, the
$1$-NN at $1$ is again a tie between $(2,1)$ (now index~0) and $(0,0)$ (now
index~1), now resolved to $(2,1)$, predicting $1$.  The loss at $(1,0)$ becomes
$1$, giving $\Delta_{\mathrm{del}}(S,0,1,0)=1$.
